\section{An exact deficit identity}\label{sec:4}

The contraction theorem gives an upper bound. To understand equality and near equality, we need to identify where the proof can lose volume. The certificate separates these losses into three terms.

\subsection{The two losses in the certificate inequality}\label{sec:4.1}

For a feasible primal--dual pair in Section~\ref{sec:2}, let
\begin{equation}\label{eq:4.1}
\tau_r=\beta_r-g_r^Tc,\qquad x_r=Ag_r,\qquad
\gamma_r=\alpha_r\tau_r-q_r^Tx_r\ge0,
\qquad Y=M^{1/2}AM^{1/2}.
\end{equation}

Each $\gamma_r$ is a cone pairing. The identities
$\operatorname{tr}Y=\operatorname{tr}(AM)$ and
$\det Y=\det A\det M$ give
\begin{equation}\label{eq:4.2}
\begin{aligned}
\mathcal C&:=\sum_r\gamma_r
=\sum_r\alpha_r\beta_r-c^T\sum_r\alpha_rg_r-\operatorname{tr}(AM)\\
&=n-\operatorname{tr}(AM),\\
\mathcal D&:=\operatorname{tr}Y-n-\log\det Y\ge0,\\
G&:=-\log\det M-\log\det A=\mathcal C+\mathcal D.
\end{aligned}
\end{equation}

Here $\mathcal C$ records the slack in the aggregated support inequalities. The matrix divergence $\mathcal D$ records the slack in the spectral inequality $\log\det Y\le\operatorname{tr}Y-n$. Its nonnegativity follows from $r-1-\log r\ge0$ for each positive eigenvalue $r$. The gap $G$ is computable from a feasible ellipsoid and a feasible certificate.

\subsection{Add the geometric loss}\label{sec:4.2}

For the explicit certificate of Section~\ref{sec:3}, define
$\Delta_{\rm geom}=\log(R_n\det M)=\log(R_n\det Q)$. Theorem~\ref{thm-frame-bound} proves that this term is nonnegative and controls $\chi$.

\begin{proposition}[{exact volume deficit}]\label{prop-gap}
In John coordinates, every feasible cut ellipsoid satisfies
\begin{equation}\label{eq:4.3}
\boxed{\epsilon:=\log\frac{R_n}{\det A}
=\Delta_{\rm geom}+\mathcal C+\mathcal D.}
\end{equation}

All three terms are nonnegative. For the actual new John ellipsoid, the same deficit is $\epsilon=\log(R_nw(K)/w(K\cap H))$ in the original coordinates.

\end{proposition}

\begin{proof}
Add the geometric loss to the primal--dual gap. The intermediate
certificate determinant cancels, leaving exactly the volume deficit:
\begin{equation}\label{eq:4.4}
\begin{aligned}
\Delta_{\rm geom}+\mathcal C+\mathcal D
&=\log(R_n\det M)-\log\det M-\log\det A\\
&=\log\frac{R_n}{\det A}=\epsilon.
\end{aligned}
\end{equation}

At the new optimum, $\det A=w(K\cap H_v)/w(K)$. Affine invariance gives the stated interpretation. 
\end{proof}

The three terms have distinct zero sets:
\begin{equation}\label{eq:4.5}
\begin{aligned}
\Delta_{\rm geom}=0&\iff\chi=0\iff s_i\in\{-1,b\}\quad\forall i,\\
\mathcal C=0&\iff\gamma_r=0\quad\forall r,\\
\mathcal D=0&\iff Y=I\iff A=M^{-1}.
\end{aligned}
\end{equation}

Thus zero deficit imposes the model contact geometry, complementarity, and shape matching simultaneously. For small deficit, the same decomposition gives quantitative control; in particular, $\chi\le1500\Delta_{\rm geom}\le1500\epsilon$.

\subsection{Small matrix divergence means close shape matching}\label{sec:4.3}

We use $\|\cdot\|_F$ for the Frobenius matrix norm; vector norms are Euclidean.

\begin{lemma}[{matrix matching}]\label{lem-matching}
For any normalized feasible pair, if $G\le1/4$, then
\begin{equation}\label{eq:4.6}
\|M^{1/2}AM^{1/2}-I\|_F\le2\sqrt G.
\end{equation}

\end{lemma}

\begin{proof}
We first use the small divergence to bound each eigenvalue from
above. On that bounded range, the same scalar divergence controls squared
distance from $1$. Put $h(t)=t-1-\log t$.
Integrating $h'(t)=1-1/t$ on either side of $1$ gives
\begin{equation}\label{eq:4.7}
h(t)\ge\frac{(t-1)^2}{2\max\{1,t\}}\quad(t>0),
\qquad h(2)=1-\log2>\frac14.
\end{equation}

For the eigenvalues $\rho_j>0$ of $Y$,
$\sum_jh(\rho_j)=\mathcal D\le G\le1/4$.
Since $h$ is increasing on $[1,\infty)$, each $\rho_j<2$. Hence
\begin{equation}\label{eq:4.8}
\|Y-I\|_F^2=\sum_j(\rho_j-1)^2
\le4\sum_jh(\rho_j)=4\mathcal D\le4G.
\end{equation}

Taking square roots proves the claim. 
\end{proof}

Section~\ref{sec:9} also gives an exact sum-of-squares identity for each cone pairing and a corresponding weighted estimate for multiplier directions. For the geometric conclusions below, the essential point is already visible: a small gap enforces matrix matching, while the geometric part of the deficit controls how close $M$ is to the model certificate.
